\documentclass[11pt,a4paper]{article}
\usepackage[margin=23mm]{geometry}
\usepackage{fontspec}
\setmainfont{Latin Modern Roman}
\setmonofont{Latin Modern Mono}[Scale=0.85]
\usepackage{amsmath,amssymb,amsthm,booktabs,xurl,hyperref}
\setlength{\parindent}{0pt}\setlength{\parskip}{6pt}
\setlength{\emergencystretch}{3em}
\newcommand{\code}[1]{\nolinkurl{#1}}
\newcommand{\E}{\mathbb E}
\newcommand{\Law}{\operatorname{Law}}
\newtheorem{theorem}{Twierdzenie}
\title{S05 / SAMPLER\_004\\\large Naprawa gałęzi $h=0$ z pełnym momentem odpowiedzi}
\author{Notatka robocza dla Niirmata}
\date{10 października 2026}
\begin{document}\maketitle

\textbf{Wynik częściowy:} skonstruowano wykonalny sampler dla $h=0$
z punktową dominacją pełnego prawa odpowiedzi, zatem z absolutną ciągłością
i kierunkowym momentem $\sum J_0^2/P_0\le1+e_0$, gdzie $e_0<2^{-50}$.
Dowód jest tekstowy, liczby sprawdzono w SageMath 10.9 przez $\mathbb Z$,
$\mathbb Q$ i balls512. Nie wykonano nowego sprawdzenia kernelowego.

\textbf{Cel dla wszystkich $h$ pozostaje otwarty.} To usunięcie konkretnej
przeszkody poprzedniego kandydata przy $h=0$, nie ograniczenie końcowego
kwantyfikatora i nie twierdzenie bezpieczeństwa FT1536.
Nie przypisujemy zerowego klucza do rozkładu emitted KeyGen.

Po decyzji właściciela cała nowa praca powstaje bezpośrednio w
\code{development/S05_QROM/} pod \code{proofs/ft1536/}; runtime jest w
\code{.build/sampler_004/}. Historyczne W zachowano, a przejście tej samej
sesji zapisuje \code{SAMPLER_004_HANDOFF.json}.

\section{Dokładne prawo docelowe}
Baza źródeł to commit \code{0a6999d7} i piny SAMPLER-003.
Używamy $q=18433$, $\sigma=768$, $B=2093922385$ oraz boxu
$[-65535,65535]^{3072}$. Nie zmieniamy modelowej normy, cap16 ani emisji.
Niech $q_2(x,y)=x^2+xy+y^2$, a $Q_0(v)$ będzie sumą 768 takich bloków.
Waga połówki to $w(v)=\exp(-Q_0(v)/(2\sigma^2))$.

W uczciwym $P_0$ wyzwanie $c\in R_q$ jest jednostajne. Jedna próba pobiera
parę $(z_1,z_2)$ z wagą $w(z_1)w(z_2)$ na włóknie
$z_1\bmod q=c$, następnie sprawdza $Q_0(z_1)+Q_0(z_2)<B$.
Normę można odrzucić najwyżej 16 razy. Po pierwszej akceptacji następuje
jedna modelowa emisja; jej porażka również daje \code{none}.

Dla $h=0$ prawo przed normą rozkłada się dokładnie na:
768 niezależnych dwuwymiarowych Gaussów w klasach reszt pierwszej połówki
oraz 768 niezależnych Gaussów na pełnym boxie drugiej połówki.
Właśnie tę faktoryzację implementujemy z kontrolą punktową.
Wiadomość, nonce i stan są ignorowane, więc wywołanie ze stałym
publicznym \code{initial} nie wymaga żadnego dostępu do historii hasza.

\section{Wykonalne losowanie we włóknie}
Dla pary reszt $(a,b)$ wypisujemy wszystkie
$(x,y)\in[-65535,65535]^2$ z $x\equiv a$, $y\equiv b\pmod q$.
Jest ich co najmniej 1 i najwyżej $K=64$.
Niech $m$ będzie najmniejszym $q_2(x,y)$ w tej liście i
$u_i=\exp(-(q_2(x_i,y_i)-m)/(2\sigma^2))$.
Co najmniej jedna waga wynosi dokładnie 1, a każda leży w $(0,1]$.
Odejmowanie $m$ nie zmienia znormalizowanego prawa.

\subsection*{Jednolita aproksymacja wykładnicza, bez nieznanego czasu pętli}
Funkcja \code{exp_weight_interval} zwraca całkowite końce przedziału
w skali $2^{384}$ dla $e^{-t}$ przy dowolnym wymiernym $t\ge0$.
Dla $t=0$ przedział jest dokładny. Dla $t\ge400$ używamy
$[0,2^{-384}]$: $e>2$ daje $e^{-t}\le2^{-400}<2^{-384}$.

Dla $0<t<400$ rozwijamy $e^{-t/512}$ naprzemiennym szeregiem Taylora.
Częściowa suma stopnia 127 jest dolną, a stopnia 128 górną granicą,
ponieważ $t/512<25/32<1$ i kolejne wyrazy maleją.
Szerokość wynosi najwyżej $1/128!<2^{-700}$.
Zaokrąglamy końce na zewnątrz do siatki $2^{-384}$ i dziewięć razy
podnosimy przedział do kwadratu, za każdym razem zaokrąglając na zewnątrz
i przecinając z $[0,1]$. W jednostkach siatki szerokość spełnia
$w_0\le3$, $w_{i+1}\le2w_i+2$, zatem $w_9\le2558<4096$.
Otrzymujemy jednolicie szerokość $\delta_w=2^{-372}$.
Wszystkie operacje wykonawcze są wymierne lub całkowite, z najwyżej
128 wyrazami i dziewięcioma kwadratami; nie ma nieograniczonego
zwiększania precyzji ani pomijanej gałęzi błędu numerycznego.

\subsection*{Dyadyczne prawo kategorii}
Niech $\ell_i\le u_i\le v_i$ będą otrzymanymi granicami, a $U=\sum_i v_i$.
Dla $N=2^{256}$ przydzielamy najpierw
$n_i=\lfloor N\ell_i/U\rfloor$ bitstringów.
Całą pozostałą masę przydzielamy wybranemu minimizerowi $q_2$,
czyli kategorii o największym prawdziwym prawdopodobieństwie.
Tak powstałe liczności sumują się dokładnie do $N$.

Dla kategorii innych niż ta wyróżniona masa nie wzrasta względem
prawdziwego prawa. Całkowity ubytek masy jest co najwyżej
$K\delta_w+K/N$, ponieważ prawdziwy normalizator jest co najmniej 1.
Prawdopodobieństwo wyróżnionej kategorii wynosi co najmniej $1/K$.
Zatem dla \emph{każdej} pary reszt i każdej kategorii
\[
 j_{\mathrm{cat}}(i)\le K_{\mathrm{cat}}p_{\mathrm{cat}}(i),\qquad
 K_{\mathrm{cat}}=1+64^2(2^{-372}+2^{-256}).
\]
To jednostronna dominacja punktowa, a nie zamiana błędu TV na moment.

\section{Druga połówka: kontrolowana masa wektora awaryjnego}
Używamy przypiętych tablic SAMPLER-003: $T=15360$, $p=256$, $b=384$,
\[
 W_x=\lceil2^p e^{-x^2/(4\sigma^2)}\rceil,\qquad
 A_s=\lfloor2^p e^{-s^2/(4\sigma^2)}\rfloor.
\]
Sposób losowania współrzędnych modulo $S=\sum_xW_x$ i akceptacji
przez $A_{x+y}$ pozostaje taki sam. Zmieniamy \emph{wewnętrzne}
wyczerpanie 192 prób pary: zwracamy parę $(0,0)$ zamiast kończyć Sign
syntetyczną porażką. Ta para ma dodatnią, policzalną masę w docelowym
Gaussie i jest następnie poddana zwykłej normie całego wektora.
Nie ukrywamy tej gałęzi; jej koszt dystrybucyjny jest składnikiem poniżej.

Niech $Z_T=\sum_{|x|\le T}e^{-x^2/(4\sigma^2)}$ i $Z_\triangle$ będzie
normalizatorem Gaussa $e^{-q_2/(2\sigma^2)}$ na pełnym boxie pary.
Dokładna akceptacja dyadycznej propozycji to przypięte $\alpha>1/2$.
Mamy
\[
 Z_T\ge(S-(2T+1))/2^p=:Z_-,\qquad
 \kappa=(1+S/2^b)(1+2^{-p}e^{T^2/(4\sigma^2)}).
\]
Każda współrzędna propozycji ma masę co najwyżej
$\kappa e^{-x^2/(4\sigma^2)}/Z_T$.
Próg akceptacji zaokrąglamy w dół, więc po warunkowaniu na akceptację
masa pary względem docelowego Gaussa jest ograniczona czynnikiem
$\kappa^2 Z_\triangle/(\alpha Z_T^2)$.
Prawdopodobieństwo wyczerpania pary to $\delta_p=(1-\alpha)^{192}$.
Dodanie go do $(0,0)$ zwiększa czynnik co najwyżej o $\delta_pZ_\triangle$.

Górną granicę $Z_+$ na $Z_\triangle$ obliczamy przez dopełnienie kwadratu
i rozdzielenie parzystych oraz nieparzystych wartości $y$:
\[
Z_\triangle\le
 \overline\theta_0(a)\overline\theta_0(3a)
 +\overline\theta_{1/2}(a)\overline\theta_{1/2}(3a)=Z_+,\qquad
 a=1/(2\sigma^2),\quad
 \theta_s(a)=\sum_{k\in\mathbb Z}e^{-a(k+s)^2}.
\]
Przez $\overline\theta_s$ oznaczamy jawną majorantę $\theta_s$:
sumy do 24576 są liczone w balls512, a oba ogony przez szereg geometryczny:
pierwszy pominięty wyraz $e^{-ar^2}$ ma następniki z ilorazem najwyżej
$e^{-a(2r+1)}$. Suma po wszystkich liczbach całkowitych jest majorantą
boxu; nie zakładamy równości boxu z nieskończoną kratą.

Ostatecznie rzeczywista para, \emph{łącznie z fallbackiem}, spełnia
\[
 j_{\mathrm{pair}}(z)\le K_{\mathrm{pair}}p_{\mathrm{pair}}(z),\qquad
 K_{\mathrm{pair}}=
 \frac{\kappa^2 Z_+}{\alpha Z_-^2}+\delta_p Z_+.
\]
Rachunek daje $K_{\mathrm{pair}}-1<2^{-108}$
(około $4.81075\cdot10^{-34}$).

\section{Pełna odpowiedź: porażki i końcowy certyfikat}
Jedna nieodrzucona jeszcze próba przy ustalonym $c$ spełnia dominację
z czynnikiem
$K_A=(K_{\mathrm{cat}}K_{\mathrm{pair}})^{768}$.
Zastosowanie tej samej normy jest deterministycznym pushforwardem,
więc zachowuje dominację dla \code{some} i \code{none}.

Cały cap16 można przedstawić jako deterministyczną funkcję 16 niezależnych
prób: zwróć pierwszy zaakceptowany wektor, a przy braku akceptacji porażkę;
następnie wykonaj jedną emisję. Dominacja produktu ma czynnik $K_A^{16}$.
Sumowanie mas po włóknach tej funkcji daje
\[
 J_0(o\mid c)\le K_A^{16}P_0(o\mid c)
 \quad\text{dla każdego }c,o,
\]
również dla wyczerpania cap16 i porażki modelowej emisji.
Fakt, że druga połówka symulatora zawsze mieści się w signed16,
nie usuwa porażek emisji z prawa docelowego ani z tego oszacowania.

Algorytm zawiera dokładny skrót: oblicza sumę minimalnych norm bloków
pierwszej połówki przy $c$. Jeśli ta suma jest co najmniej $B$, to wszystkie
wektory w tym włóknie oblewają normę niezależnie od drugiej połówki.
Uczciwe prawo wtedy daje \code{none} z prawdopodobieństwem 1, więc bezpośredni
zwrot porażki jest dokładną realizacją tej gałęzi, nie przybliżeniem.

Wyzwanie losujemy przez 1536 niezależnych 80-bitowych słów modulo $q$.
Każda masa współczynnika jest co najwyżej $(1+q/2^{80})/q$, stąd
$K_H=(1+q/2^{80})^{1536}$ jest czynnikiem względem idealnego jednostajnego
wyzwania. Dokładnej jednostajności nie zakładamy.

\begin{theorem}[Zakres $h=0$, idealne bity, model współczynnikowy]
Opisany wykonywalny sampler pełnej pary $(c,o)$ spełnia
\[
 J_0(c,o)\le C P_0(c,o),\qquad C=K_H K_A^{16}.
\]
W konsekwencji $J_0\ll P_0$ oraz
\[
 \sum_{P_0(c,o)>0}\frac{J_0(c,o)^2}{P_0(c,o)}\le C^2=1+e_0,
 \qquad 0<e_0<2^{-50}.
\]
\end{theorem}
\begin{proof}
Pierwsza nierówność wynika z przemnożenia wykazanych czynników wyzwania
i odpowiedzi. Gdy $P_0(c,o)=0$, dominacja i nieujemność dają $J_0(c,o)=0$.
Dla dodatniej masy mianownika $J_0^2/P_0\le C^2P_0$;
sumowanie i normalizacja $P_0$ dają wniosek.
Nierówność liczbowa została sprawdzona w balls512 na dokładnych
parametrach i przypiętych tablicach; przedział dla $e_0$ ma środek
około $4.6840074950\cdot10^{-17}$.
\end{proof}

Zachowano konserwatywną postać $C^2$, odpowiadającą dotychczasowemu
narzędziu punktowemu; nie podmieniono po cichu kierunku $J/P$ ani funkcji
straty. Jest to certyfikat specjalizacji przy $h=0$, a nie instancja
całego typu \code{LocalJointCertificate}, który kwantyfikuje po wszystkich $h$.

\section{Dlaczego faktoryzacja nie kończy ogólnego przypadku}
Dla dowolnego $h$ oznaczmy
\[
 Z_1(r)=\sum_{z_1:\,z_1\bmod q=r}w(z_1),\qquad Z_G=\sum_{z_2}w(z_2),
 \qquad r(z_2)=c-hz_2\pmod q.
\]
Naturalna próba rozszerzenia najpierw losuje $z_2$ z prawa $U=w/Z_G$,
a następnie $z_1$ z Gaussa warunkowanego na resztę $r(z_2)$.
Nie jest to w ogólności uczciwy Gaussian na całym włóknie.
Jego gęstość to $J=w(z_1)w(z_2)/(Z_G Z_1(r(z_2)))$,
podczas gdy $P=w(z_1)w(z_2)/Z_{h,c}$, gdzie
$Z_{h,c}=\sum_{z_2}w(z_2)Z_1(r(z_2))$.
Z bezpośredniego sumowania po $z_1$ dostajemy dokładnie
\[
 \sum\frac{J^2}{P}=
 \E_{z_2\sim U}[Z_1(r(z_2))]\,
 \E_{z_2\sim U}[1/Z_1(r(z_2))].
\]
Wszystkie normalizatory są dodatnie: każda klasa reszt ma reprezentanta
w pierwszym boxie. Dla $h=0$ normalizator $Z_1(c)$ jest stały i iloczyn
wynosi 1. Dla innych $h$ trzeba kontrolować \emph{także odwrotność}
normalizatora. Sama kontrola średniej lub równość nośnika nie wystarcza.
To tożsamość dla prawa przed bramką normy, nie gotowy certyfikat całego Sign.

\code{RESIDUAL_NORMALIZER.sage} sprawdził tożsamość dokładnie w 34 małych
modelach wymiernych dla $q=3,5$. To test implementacji sumowania;
ogólny argument jest powyższym rachunkiem, a nie ekstrapolacją testów.
Nie ogłaszamy niemożliwości innej konstrukcji dla pozostałych kluczy.

\section{Wykonanie, koszt i status}
Źródło: \code{HZERO_SAMPLER.sage}. Rachunek główny i powtórzenie po
przejściu do development zakończyły się poprawnie. Sprawdzono:
sukces dla wyzwania zerowego, sukces przy granicy wyboru reprezentanta,
pełne wyczerpanie cap16 oraz dokładny abort z dolnej granicy normy.
Wymuszono również fallback pary przy rzeczywistym limicie 192.
Każde zdarzenie ma swój wpis w receipcie; nie warunkujemy gry na sukces.

Losowość testowa jest deterministycznym SHAKE z publicznymi etykietami.
Dowód rozkładu zakłada idealne niezależne bity i nie dowodzi bezpieczeństwa PRG.
Ograniczenia na odpowiedź:
\begin{itemize}
\item 49152 przedziały wag we włóknach; każdy wymaga najwyżej 128 wyrazów
Taylora i dziewięciu kwadratów, bez nieograniczonego retry numerycznego;
\item 2359296 propozycji par drugiej połówki;
\item 2419187712 bitów, czytanych leniwie, bez alokowania całej taśmy;
\item dodatkowo deterministyczne wyliczenie klas, CDF, norm i odczyt
przypiętych tablic. To policzone operacje algorytmu, nie pełny certyfikat
maszyny bitowej ani benchmark produkcyjny.
\end{itemize}

Nie ponawiano niezmienionej próby Lean, której sam import poprzednio
przekroczył limit. Nowy dowód pozostaje tekstowy i wymaga niezależnego
odbioru oraz późniejszej formalizacji. Nie zmieniono T12.1/B20, C,
paperu ani strony; zapis Git jest lokalny, bez push.

\textbf{Co dalej:} trzeba rozwiązać problem rozkładu syndromów lub
normalizatora reszt dla pozostałych $h$, utrzymując pełne prawo porażek.
Wynik SAMPLER-004 pozwala zachować gałąź $h=0$ w przyszłym selektorze
publicznym, ale nie dowodzi niczego o momencie jego innych gałęzi.
Cel końcowy nadal ma pierwotny zakres wszystkich $h$.
\end{document}
